\begin{corollary}[Archimedean Reciprocal Form]
\label{cor:archimedean-reciprocal-form}
For every $\varepsilon\in\mathbb{R}$ with $\varepsilon>0$, there exists
$n\in\mathbb{N}$ such that
\[
  0<\frac{1}{n}<\varepsilon.
\]
Equivalently,
\[
  (\forall\varepsilon>0)(\exists n\in\mathbb{N})\left(\frac{1}{n}<\varepsilon\right).
\]

\smallskip
\noindent
\hyperref[prf:archimedean-reciprocal-form]{\textit{Go to proof.}}
\end{corollary}
\end{corollarybox}

\begin{remark*}[Standard quantified statement]
\[
\forall \varepsilon\in\mathbb{R}\;
\bigl(\varepsilon>0\Longrightarrow
\exists n\in\mathbb{N}\;(0<1/n<\varepsilon)\bigr).
\]
\end{remark*}

\begin{remark*}[Predicate reading]
\[
P_{\mathbb{R}}=\mathsf{OrderedSet}(\mathbb{R},\leq),\qquad
\operatorname{HasArchimedeanProperty}(P_{\mathbb{R}})
\Longrightarrow
\forall \varepsilon>0\;\exists n\in\mathbb{N}\;(0<1/n<\varepsilon).
\]
\end{remark*}

\begin{remark*}[Negated quantified statement]
\[
\exists \varepsilon>0\;\forall n\in\mathbb{N}\;\left(\frac{1}{n}\ge \varepsilon\right).
\]
\end{remark*}

\begin{remark*}[Negation predicate reading]
\[
\exists \varepsilon>0\;\forall n\in\mathbb{N}\;(1/n\ge \varepsilon).
\]
\end{remark*}

\begin{remark*}[Failure modes]
\begin{description}
  \item[Exposition.]
The corollary fails only if some positive scale remains below every reciprocal
natural number.


  \item[\exists-condition.]
  \[
\exists \varepsilon>0\;\forall n\in\mathbb{N}\;\left(\frac{1}{n}\ge \varepsilon\right).
\]
\[
\exists \varepsilon>0\;\forall n\in\mathbb{N}\;(1/n\ge \varepsilon).
\]
\end{description}
\end{remark*}

\begin{remark*}[Failure modes]
\begin{description}
  \item[Exposition.]
\[
\underbrace{\varepsilon>0}_{\text{positive scale}}
\qquad\text{and}\qquad
\underbrace{\forall n\in\mathbb{N}\;(1/n\ge\varepsilon)}_{\text{reciprocals never pass below it}}.
\]


  \item[\exists-condition.]
  \[
\exists \varepsilon>0\;\forall n\in\mathbb{N}\;\left(\frac{1}{n}\ge \varepsilon\right).
\]
\[
\exists \varepsilon>0\;\forall n\in\mathbb{N}\;(1/n\ge \varepsilon).
\]
\end{description}
\end{remark*}

\begin{remark*}[Interpretation]
The reciprocals of the natural numbers descend below every positive real
tolerance.
\end{remark*}

\begin{dependencies}
\begin{itemize}
\item \hyperref[def:reals]{Real Numbers}.
\item \hyperref[def:natural-numbers]{Natural Numbers}.
\item \hyperref[def:ordered-set]{Ordered Set}.
\item \hyperref[thm:archimedean-property]{Archimedean Property}.
\end{itemize}
\end{dependencies}



